\section{Contexto institucional}\label{sec:institutional}

\subsection{El salario mínimo peruano}

El Perú fija un único salario mínimo nacional, la \emph{Remuneración Mínima Vital} (RMV), que se aplica a los trabajadores del régimen laboral de la actividad privada. A diferencia de sistemas federales como los de Estados Unidos o Brasil, y de los países con pisos salariales por ocupación o por región, el Perú no tiene un mínimo legal subnacional ni específico por industria en el régimen general. Un pequeño número de regímenes especiales está indexado a la RMV: la minería tiene una bonificación legal sobre ella, el régimen agrario de la Ley 31110 fija una remuneración diaria anclada en ella y el régimen de la micro y pequeña empresa (MYPE) aplica el mismo piso. Ninguno de ellos introduce variación transversal en el nivel de la política.\footnote{Como el régimen agrario paga una remuneración diaria bajo reglas propias, identifico la agricultura en todo el análisis y muestro la robustez de los resultados a su exclusión (Sección~\ref{sec:robustness}).} La consecuencia para el diseño de investigación es decisiva: como el nivel de la RMV es común a todo el país, no existe variación transversal en la política ni un calendario escalonado de adopción. Cualquier diseño de investigación creíble debe, por tanto, aprovechar las diferencias en el grado de exposición de los trabajadores a un cambio que llega a todos al mismo tiempo.

En principio, la RMV debe ser ajustada periódicamente por el Consejo Nacional de Trabajo, de carácter tripartito, en función de la inflación y la productividad. En la práctica, el Consejo nunca ha fijado la RMV, y su nivel lo determina discrecionalmente el Poder Ejecutivo mediante un Decreto Supremo. En consecuencia, los ajustes son poco frecuentes e irregulares tanto en su calendario como en su magnitud, y el Perú figura entre los países de la región que fijan su salario mínimo de forma más volátil \citep{jimenez_2023, cespedes_2005}. La RMV se había mantenido en 930 soles desde abril de 2018.\footnote{Fijada por el Decreto Supremo 004-2018-TR. Los dos aumentos anteriores elevaron la RMV de 750 a 850 soles en mayo de 2016 y de 850 a 930 soles en abril de 2018, de modo que la reforma de 2022 siguió a un congelamiento inusualmente largo.}

Dos características hacen que la RMV sea inusualmente vinculante para los estándares internacionales. Primero, su relación con el salario mediano es alta. En mi muestra previa a la reforma, el salario mensual equivalente mediano nacional de un asalariado privado es de 1000 soles, de modo que el mínimo anterior de 930 equivalía al 93 por ciento de la mediana, el nuevo mínimo quedó ligeramente por encima de ella, y el índice de Kaitz, la razón entre el mínimo y el salario mediano, es cercano a uno a nivel nacional y se aproxima a dos en los departamentos más pobres. Segundo, el cumplimiento dista de ser universal. En línea con trabajos anteriores \citep{jaramillo_lopez_2006, cespedes_2005}, una gran proporción de trabajadores percibe menos que el piso legal, especialmente en las zonas rurales, en las empresas pequeñas y en el sector informal, mientras que el cumplimiento dentro de las empresas formales registradas es alto. La capacidad de fiscalización es limitada.\footnote{La inspección laboral está a cargo de la Superintendencia Nacional de Fiscalización Laboral (SUNAFIL), cuyo personal inspector es escaso fuera de Lima. Trabajos anteriores estiman el incumplimiento total de la RMV en alrededor de 30 a 40 por ciento, concentrado en las zonas rurales, las empresas pequeñas y el sector informal \citep{jaramillo_lopez_2006, cespedes_2005}.} Estas características implican que un aumento nominal de la RMV puede ser mecánicamente vinculante para los trabajadores formales cercanos al piso y, a la vez, dejar a buena parte del sector informal sin cobertura formal, de modo que el alcance real de la política depende del grado de informalidad y de la facilidad de movimiento entre los márgenes formal e informal.

\subsection{La reforma de 2022}

La reforma que estudio fue dispuesta por el Decreto Supremo 003-2022-TR, publicado el 3 de abril de 2022, y entró en vigor el 1 de mayo de 2022. Elevó la RMV de 930 a 1025 soles, un aumento de 95 soles o 10.2 por ciento, el primer cambio en cuatro años. La Figura~\ref{fig:mw} muestra el valor nominal y real de la RMV durante 2021--2025. Conviene destacar dos puntos. El aumento nominal de 10.2 por ciento fue compensado en parte por la inflación, que rondó el 8 por ciento durante 2022, de modo que el aumento del valor real del piso fue más modesto y se erosionó en los dos años siguientes. Esto modera el tamaño esperado de cualquier respuesta salarial y es relevante al comparar mis estimaciones con las de episodios de menor inflación. La reforma ocurrió, además, durante la recuperación posterior a la pandemia, un período de turbulencia inusual en el mercado laboral que abordo de forma explícita en el análisis empírico.

Una segunda reforma, dispuesta por el Decreto Supremo 006-2024-TR y vigente desde el 1 de enero de 2025, elevó la RMV a 1130 soles. Este cambio queda justo fuera de mi ventana principal de estudio, pero los trimestres de la ENAHO de 2025 están disponibles y utilizo esta segunda reforma como una réplica fuera de muestra del diseño en la Sección~\ref{sec:robustness}. Ambas reformas son de magnitud muy similar, 10.2 por ciento cada una, lo que convierte al episodio de 2025 en un ejercicio natural de validación en un entorno más tranquilo, posterior a la pandemia.

\begin{figure}[!t]
\centering
\includegraphics[width=0.82\textwidth]{fig1_minimum_wage.pdf}
\caption{Salario mínimo peruano nominal y real, 2021--2025. La serie real está deflactada por el índice de precios al consumidor de Lima (diciembre de 2021 = 100). Las líneas verticales punteadas marcan las reformas vigentes desde el 1 de mayo de 2022 (Decreto Supremo 003-2022-TR) y el 1 de enero de 2025 (Decreto Supremo 006-2024-TR).}
\label{fig:mw}
\end{figure}
